% =====================================================================
%  Capítulo 5 — Desenvolvimento (arquitetura e implementação)
% =====================================================================
\chapter{Desenvolvimento}
\label{cap:desenvolvimento}

Este capítulo descreve a arquitetura e a implementação do ARGUS, relacionando
cada módulo aos requisitos que implementa. A descrição é fiel ao código
efetivamente produzido e testado, cujos artefatos principais residem em
`backend/` e `frontend/`.

\section{Visão geral da arquitetura}

O sistema é organizado em três camadas: \textbf{front-end} (aplica\c{c}\~ao web),
\textbf{back-end} (API e l\'ogica de dom\'inio) e \textbf{banco de dados}
(PostgreSQL). A Figura~\ref{fig:arquitetura} apresenta a arquitetura de
m\'odulos e os fluxos principais.

\begin{figure}[htbp]
\centering
% Diagrama: Arquitetura do ARGUS (visão de contêineres/módulos)
% Uso: % Diagrama: Arquitetura do ARGUS (visão de contêineres/módulos)
% Uso: % Diagrama: Arquitetura do ARGUS (visão de contêineres/módulos)
% Uso: \input{diagrams/arquitetura} dentro de um ambiente figure.
\begin{tikzpicture}[
  font=\small,
  fe/.style={rectangle, draw=black!60, fill=blue!6, rounded corners=2pt, align=center, inner sep=5pt, minimum width=3.8cm, text width=3.6cm},
  be/.style={rectangle, draw=black!60, fill=green!6, rounded corners=2pt, align=center, inner sep=5pt, minimum width=3.8cm, text width=3.6cm},
  db/.style={rectangle, draw=black!60, fill=orange!8, rounded corners=2pt, align=center, inner sep=5pt, minimum width=3.8cm, text width=3.6cm},
  lab/.style={font=\scriptsize},
  arrow/.style={-{Stealth[length=2.4mm]}, semithick}
]
  % Coluna front-end
  \node[fe] (ui) at (0,2.4) {Painel de simula\c{c}\~ao};
  \node[fe] (fftchart) at (0,0) {Gr\'afico de FFT + limiar\\ e ordens $N$};
  \node[fe] (check) at (0,-2.4) {Checklist de valida\c{c}\~ao};

  % Coluna back-end
  \node[be] (api) at (7,2.4) {API REST\\ \scriptsize POST /simulacoes, /snapshots};
  \node[be] (gen) at (7,0) {Gera\c{c}\~ao de sinal};
  \node[be] (fftm) at (7,-2.4) {FFT e extra\c{c}\~ao $R^3$};
  \node[be] (disc) at (7,-4.8) {Motor de descarte\\ \scriptsize regra de ouro + Paralelep\'ipedo};
  \node[be] (pers) at (7,-7.2) {Persist\^encia};

  % Coluna banco
  \node[db] (hier) at (14.5,3.6) {Planta/Área/\\M\'aquina/Ponto};
  \node[db] (snap) at (14.5,1.2) {snapshots\_defeito};
  \node[db] (leit) at (14.5,-1.2) {leituras\_persistidas};
  \node[db] (trash) at (14.5,-3.6) {leituras\_trash};

  % Fluxo front -> back
  \draw[arrow] (ui.east) -- node[lab,above]{POST /simulacoes} (api.west);
  \draw[arrow] (check.east) -- (api.west);
  \draw[arrow] (api.west) to[out=180,in=0] node[lab,above,near end]{resposta (sinal+FFT+RMS)} (fftchart.east);

  % Fluxo interno do back-end
  \draw[arrow] (api) -- (gen);
  \draw[arrow] (gen) -- (fftm);
  \draw[arrow] (fftm) -- (disc);
  \draw[arrow] (disc) -- (pers);

  % Fluxo back -> banco
  \draw[arrow] (api.east) to[out=0,in=180] (snap.west);
  \draw[arrow] (pers.east) -- node[lab,above,pos=0.4]{aprovada} (leit.west);
  \draw[arrow] (pers.east) -- node[lab,below,pos=0.3]{descartada} (trash.west);
  \draw[arrow] (pers.east) to[out=0,in=180] node[lab,above,pos=0.6]{l\^e/atualiza} (hier.west);
\end{tikzpicture}
 dentro de um ambiente figure.
\begin{tikzpicture}[
  font=\small,
  fe/.style={rectangle, draw=black!60, fill=blue!6, rounded corners=2pt, align=center, inner sep=5pt, minimum width=3.8cm, text width=3.6cm},
  be/.style={rectangle, draw=black!60, fill=green!6, rounded corners=2pt, align=center, inner sep=5pt, minimum width=3.8cm, text width=3.6cm},
  db/.style={rectangle, draw=black!60, fill=orange!8, rounded corners=2pt, align=center, inner sep=5pt, minimum width=3.8cm, text width=3.6cm},
  lab/.style={font=\scriptsize},
  arrow/.style={-{Stealth[length=2.4mm]}, semithick}
]
  % Coluna front-end
  \node[fe] (ui) at (0,2.4) {Painel de simula\c{c}\~ao};
  \node[fe] (fftchart) at (0,0) {Gr\'afico de FFT + limiar\\ e ordens $N$};
  \node[fe] (check) at (0,-2.4) {Checklist de valida\c{c}\~ao};

  % Coluna back-end
  \node[be] (api) at (7,2.4) {API REST\\ \scriptsize POST /simulacoes, /snapshots};
  \node[be] (gen) at (7,0) {Gera\c{c}\~ao de sinal};
  \node[be] (fftm) at (7,-2.4) {FFT e extra\c{c}\~ao $R^3$};
  \node[be] (disc) at (7,-4.8) {Motor de descarte\\ \scriptsize regra de ouro + Paralelep\'ipedo};
  \node[be] (pers) at (7,-7.2) {Persist\^encia};

  % Coluna banco
  \node[db] (hier) at (14.5,3.6) {Planta/Área/\\M\'aquina/Ponto};
  \node[db] (snap) at (14.5,1.2) {snapshots\_defeito};
  \node[db] (leit) at (14.5,-1.2) {leituras\_persistidas};
  \node[db] (trash) at (14.5,-3.6) {leituras\_trash};

  % Fluxo front -> back
  \draw[arrow] (ui.east) -- node[lab,above]{POST /simulacoes} (api.west);
  \draw[arrow] (check.east) -- (api.west);
  \draw[arrow] (api.west) to[out=180,in=0] node[lab,above,near end]{resposta (sinal+FFT+RMS)} (fftchart.east);

  % Fluxo interno do back-end
  \draw[arrow] (api) -- (gen);
  \draw[arrow] (gen) -- (fftm);
  \draw[arrow] (fftm) -- (disc);
  \draw[arrow] (disc) -- (pers);

  % Fluxo back -> banco
  \draw[arrow] (api.east) to[out=0,in=180] (snap.west);
  \draw[arrow] (pers.east) -- node[lab,above,pos=0.4]{aprovada} (leit.west);
  \draw[arrow] (pers.east) -- node[lab,below,pos=0.3]{descartada} (trash.west);
  \draw[arrow] (pers.east) to[out=0,in=180] node[lab,above,pos=0.6]{l\^e/atualiza} (hier.west);
\end{tikzpicture}
 dentro de um ambiente figure.
\begin{tikzpicture}[
  font=\small,
  fe/.style={rectangle, draw=black!60, fill=blue!6, rounded corners=2pt, align=center, inner sep=5pt, minimum width=3.8cm, text width=3.6cm},
  be/.style={rectangle, draw=black!60, fill=green!6, rounded corners=2pt, align=center, inner sep=5pt, minimum width=3.8cm, text width=3.6cm},
  db/.style={rectangle, draw=black!60, fill=orange!8, rounded corners=2pt, align=center, inner sep=5pt, minimum width=3.8cm, text width=3.6cm},
  lab/.style={font=\scriptsize},
  arrow/.style={-{Stealth[length=2.4mm]}, semithick}
]
  % Coluna front-end
  \node[fe] (ui) at (0,2.4) {Painel de simula\c{c}\~ao};
  \node[fe] (fftchart) at (0,0) {Gr\'afico de FFT + limiar\\ e ordens $N$};
  \node[fe] (check) at (0,-2.4) {Checklist de valida\c{c}\~ao};

  % Coluna back-end
  \node[be] (api) at (7,2.4) {API REST\\ \scriptsize POST /simulacoes, /snapshots};
  \node[be] (gen) at (7,0) {Gera\c{c}\~ao de sinal};
  \node[be] (fftm) at (7,-2.4) {FFT e extra\c{c}\~ao $R^3$};
  \node[be] (disc) at (7,-4.8) {Motor de descarte\\ \scriptsize regra de ouro + Paralelep\'ipedo};
  \node[be] (pers) at (7,-7.2) {Persist\^encia};

  % Coluna banco
  \node[db] (hier) at (14.5,3.6) {Planta/Área/\\M\'aquina/Ponto};
  \node[db] (snap) at (14.5,1.2) {snapshots\_defeito};
  \node[db] (leit) at (14.5,-1.2) {leituras\_persistidas};
  \node[db] (trash) at (14.5,-3.6) {leituras\_trash};

  % Fluxo front -> back
  \draw[arrow] (ui.east) -- node[lab,above]{POST /simulacoes} (api.west);
  \draw[arrow] (check.east) -- (api.west);
  \draw[arrow] (api.west) to[out=180,in=0] node[lab,above,near end]{resposta (sinal+FFT+RMS)} (fftchart.east);

  % Fluxo interno do back-end
  \draw[arrow] (api) -- (gen);
  \draw[arrow] (gen) -- (fftm);
  \draw[arrow] (fftm) -- (disc);
  \draw[arrow] (disc) -- (pers);

  % Fluxo back -> banco
  \draw[arrow] (api.east) to[out=0,in=180] (snap.west);
  \draw[arrow] (pers.east) -- node[lab,above,pos=0.4]{aprovada} (leit.west);
  \draw[arrow] (pers.east) -- node[lab,below,pos=0.3]{descartada} (trash.west);
  \draw[arrow] (pers.east) to[out=0,in=180] node[lab,above,pos=0.6]{l\^e/atualiza} (hier.west);
\end{tikzpicture}

\caption{Arquitetura do ARGUS: front-end (React), back-end (FastAPI) e banco de dados (PostgreSQL), com fluxos de simulação, descarte e persistência.}
\label{fig:arquitetura}
\caption*{\footnotesize Fonte: elaboração própria, com base na arquitetura real do projeto.}
\end{figure}

O fluxo nominal de uma simulação é: (i) a interface envia os parâmetros a
`POST /simulacoes`; (ii) o back-end valida a entrada (incluindo o critério de
Nyquist), gera o sinal, calcula a FFT e extrai os picos \(R^3\) e as métricas;
(iii) o motor de descarte decide entre persistir e descartar com base na
última leitura persistida do ponto; (iv) a resposta retorna sinal decimado,
picos, métricas, decisão e indicadores. A Figura~\ref{fig:fluxo-descarte}
detalha o fluxo de decisão de descarte e a Figura~\ref{fig:estados} a máquina
de estados por leitura.

\begin{figure}[htbp]
\centering
% Diagrama: Fluxo de decisão do motor de descarte (regra de ouro + R^3)
% Uso: % Diagrama: Fluxo de decisão do motor de descarte (regra de ouro + R^3)
% Uso: % Diagrama: Fluxo de decisão do motor de descarte (regra de ouro + R^3)
% Uso: \input{diagrams/fluxo-descarte} dentro de um ambiente figure.
\begin{tikzpicture}[
  font=\small,
  dec/.style={diamond, draw=black!60, fill=gray!8, align=center, inner sep=1pt, aspect=3, text width=2.2cm},
  act/.style={rectangle, draw=black!60, fill=green!6, rounded corners=2pt, align=center, inner sep=5pt, minimum width=4.0cm, text width=3.8cm},
  term/.style={rectangle, draw=black!60, fill=orange!12, rounded corners=2pt, align=center, inner sep=5pt, minimum width=4.2cm, text width=4.0cm},
  arrow/.style={-{Stealth[length=2.4mm]}, semithick},
  lab/.style={font=\scriptsize}
]
  \node[act] (inicio) at (0,3.0) {Leitura avaliada (rota\c{c}\~ao + picos $R^3$)};
  \node[dec] (primeira) at (0,0.6) {Primeira\\ leitura do\\ ponto?};
  \node[dec] (ouro) at (6.0,0.6) {$|\Delta RPM| > \Delta_{rot}$\\ e n\'umero de\\ picos $>0$?};
  \node[dec] (r3) at (12.0,0.6) {Todos os picos\\ dentro do\\ Paralelep\'ipedo?};
  \node[term] (persiste) at (6.0,-2.8) {PERSISTIR\\ \scriptsize leituras\_persistidas (nova refer\^encia)};
  \node[term] (descarta) at (12.0,-2.8) {DESCARTAR\\ \scriptsize registrar em leituras\_trash};

  \draw[arrow] (inicio) -- (primeira);

  \draw[arrow] (primeira.east) -- node[lab,above]{n\~ao} (ouro.west);
  \draw[arrow] (primeira.south) -- ++(0,-1.6) -| node[lab,above,pos=0.25]{sim} (persiste.west);

  \draw[arrow] (ouro.east) -- node[lab,above]{n\~ao} (r3.west);
  \draw[arrow] (ouro.south) -- node[lab,left]{sim} (persiste.north);

  \draw[arrow] (r3.south) -- node[lab,right]{sim} (descarta.north);
  \draw[arrow] (r3.west) to[out=180,in=0] node[lab,above,pos=0.3]{n\~ao} (persiste.east);
\end{tikzpicture}
 dentro de um ambiente figure.
\begin{tikzpicture}[
  font=\small,
  dec/.style={diamond, draw=black!60, fill=gray!8, align=center, inner sep=1pt, aspect=3, text width=2.2cm},
  act/.style={rectangle, draw=black!60, fill=green!6, rounded corners=2pt, align=center, inner sep=5pt, minimum width=4.0cm, text width=3.8cm},
  term/.style={rectangle, draw=black!60, fill=orange!12, rounded corners=2pt, align=center, inner sep=5pt, minimum width=4.2cm, text width=4.0cm},
  arrow/.style={-{Stealth[length=2.4mm]}, semithick},
  lab/.style={font=\scriptsize}
]
  \node[act] (inicio) at (0,3.0) {Leitura avaliada (rota\c{c}\~ao + picos $R^3$)};
  \node[dec] (primeira) at (0,0.6) {Primeira\\ leitura do\\ ponto?};
  \node[dec] (ouro) at (6.0,0.6) {$|\Delta RPM| > \Delta_{rot}$\\ e n\'umero de\\ picos $>0$?};
  \node[dec] (r3) at (12.0,0.6) {Todos os picos\\ dentro do\\ Paralelep\'ipedo?};
  \node[term] (persiste) at (6.0,-2.8) {PERSISTIR\\ \scriptsize leituras\_persistidas (nova refer\^encia)};
  \node[term] (descarta) at (12.0,-2.8) {DESCARTAR\\ \scriptsize registrar em leituras\_trash};

  \draw[arrow] (inicio) -- (primeira);

  \draw[arrow] (primeira.east) -- node[lab,above]{n\~ao} (ouro.west);
  \draw[arrow] (primeira.south) -- ++(0,-1.6) -| node[lab,above,pos=0.25]{sim} (persiste.west);

  \draw[arrow] (ouro.east) -- node[lab,above]{n\~ao} (r3.west);
  \draw[arrow] (ouro.south) -- node[lab,left]{sim} (persiste.north);

  \draw[arrow] (r3.south) -- node[lab,right]{sim} (descarta.north);
  \draw[arrow] (r3.west) to[out=180,in=0] node[lab,above,pos=0.3]{n\~ao} (persiste.east);
\end{tikzpicture}
 dentro de um ambiente figure.
\begin{tikzpicture}[
  font=\small,
  dec/.style={diamond, draw=black!60, fill=gray!8, align=center, inner sep=1pt, aspect=3, text width=2.2cm},
  act/.style={rectangle, draw=black!60, fill=green!6, rounded corners=2pt, align=center, inner sep=5pt, minimum width=4.0cm, text width=3.8cm},
  term/.style={rectangle, draw=black!60, fill=orange!12, rounded corners=2pt, align=center, inner sep=5pt, minimum width=4.2cm, text width=4.0cm},
  arrow/.style={-{Stealth[length=2.4mm]}, semithick},
  lab/.style={font=\scriptsize}
]
  \node[act] (inicio) at (0,3.0) {Leitura avaliada (rota\c{c}\~ao + picos $R^3$)};
  \node[dec] (primeira) at (0,0.6) {Primeira\\ leitura do\\ ponto?};
  \node[dec] (ouro) at (6.0,0.6) {$|\Delta RPM| > \Delta_{rot}$\\ e n\'umero de\\ picos $>0$?};
  \node[dec] (r3) at (12.0,0.6) {Todos os picos\\ dentro do\\ Paralelep\'ipedo?};
  \node[term] (persiste) at (6.0,-2.8) {PERSISTIR\\ \scriptsize leituras\_persistidas (nova refer\^encia)};
  \node[term] (descarta) at (12.0,-2.8) {DESCARTAR\\ \scriptsize registrar em leituras\_trash};

  \draw[arrow] (inicio) -- (primeira);

  \draw[arrow] (primeira.east) -- node[lab,above]{n\~ao} (ouro.west);
  \draw[arrow] (primeira.south) -- ++(0,-1.6) -| node[lab,above,pos=0.25]{sim} (persiste.west);

  \draw[arrow] (ouro.east) -- node[lab,above]{n\~ao} (r3.west);
  \draw[arrow] (ouro.south) -- node[lab,left]{sim} (persiste.north);

  \draw[arrow] (r3.south) -- node[lab,right]{sim} (descarta.north);
  \draw[arrow] (r3.west) to[out=180,in=0] node[lab,above,pos=0.3]{n\~ao} (persiste.east);
\end{tikzpicture}

\caption{Fluxo de decisão do motor de descarte: regra de ouro e Paralelepípedo de Descarte \(R^3\).}
\label{fig:fluxo-descarte}
\caption*{\footnotesize Fonte: elaboração própria, conforme `discard_engine.py`.}
\end{figure}

\begin{figure}[htbp]
\centering
% Diagrama: Máquina de estados por leitura simulada
% Uso: % Diagrama: Máquina de estados por leitura simulada
% Uso: % Diagrama: Máquina de estados por leitura simulada
% Uso: \input{diagrams/estados} dentro de um ambiente figure.
\begin{tikzpicture}[
  font=\small,
  st/.style={rectangle, draw=black!60, fill=blue!6, rounded corners=3pt, align=center, inner sep=6pt, minimum width=3.0cm},
  fin/.style={rectangle, draw=black!60, fill=orange!12, rounded corners=3pt, align=center, inner sep=6pt, minimum width=3.0cm},
  arrow/.style={-{Stealth[length=2.4mm]}, semithick},
  lab/.style={font=\scriptsize}
]
  \node[st] (conf) at (0,0) {Configurado};
  \node[st] (sinal) at (4.0,0) {Sinal gerado};
  \node[st] (fft) at (8.0,0) {FFT/RMS/DC};
  \node[st] (aval) at (12.0,0) {Avaliando descarte};
  \node[fin] (pers) at (16.4,0) {Persistida};
  \node[fin] (desc) at (14.2,-2.4) {Descartada};

  \draw[arrow] (conf) -- node[lab,above]{gerar sinal} (sinal);
  \draw[arrow] (sinal) -- node[lab,above]{calcular FFT} (fft);
  \draw[arrow] (fft) -- node[lab,above]{avaliar} (aval);
  \draw[arrow] (aval.east) -- node[lab,above,pos=0.4]{regra de ouro / fora de $R^3$ / primeira leitura} (pers.west);
  \draw[arrow] (aval.south) to[out=-90,in=180] node[lab,left,pos=0.25]{todos dentro de $R^3$} (desc.west);
\end{tikzpicture}
 dentro de um ambiente figure.
\begin{tikzpicture}[
  font=\small,
  st/.style={rectangle, draw=black!60, fill=blue!6, rounded corners=3pt, align=center, inner sep=6pt, minimum width=3.0cm},
  fin/.style={rectangle, draw=black!60, fill=orange!12, rounded corners=3pt, align=center, inner sep=6pt, minimum width=3.0cm},
  arrow/.style={-{Stealth[length=2.4mm]}, semithick},
  lab/.style={font=\scriptsize}
]
  \node[st] (conf) at (0,0) {Configurado};
  \node[st] (sinal) at (4.0,0) {Sinal gerado};
  \node[st] (fft) at (8.0,0) {FFT/RMS/DC};
  \node[st] (aval) at (12.0,0) {Avaliando descarte};
  \node[fin] (pers) at (16.4,0) {Persistida};
  \node[fin] (desc) at (14.2,-2.4) {Descartada};

  \draw[arrow] (conf) -- node[lab,above]{gerar sinal} (sinal);
  \draw[arrow] (sinal) -- node[lab,above]{calcular FFT} (fft);
  \draw[arrow] (fft) -- node[lab,above]{avaliar} (aval);
  \draw[arrow] (aval.east) -- node[lab,above,pos=0.4]{regra de ouro / fora de $R^3$ / primeira leitura} (pers.west);
  \draw[arrow] (aval.south) to[out=-90,in=180] node[lab,left,pos=0.25]{todos dentro de $R^3$} (desc.west);
\end{tikzpicture}
 dentro de um ambiente figure.
\begin{tikzpicture}[
  font=\small,
  st/.style={rectangle, draw=black!60, fill=blue!6, rounded corners=3pt, align=center, inner sep=6pt, minimum width=3.0cm},
  fin/.style={rectangle, draw=black!60, fill=orange!12, rounded corners=3pt, align=center, inner sep=6pt, minimum width=3.0cm},
  arrow/.style={-{Stealth[length=2.4mm]}, semithick},
  lab/.style={font=\scriptsize}
]
  \node[st] (conf) at (0,0) {Configurado};
  \node[st] (sinal) at (4.0,0) {Sinal gerado};
  \node[st] (fft) at (8.0,0) {FFT/RMS/DC};
  \node[st] (aval) at (12.0,0) {Avaliando descarte};
  \node[fin] (pers) at (16.4,0) {Persistida};
  \node[fin] (desc) at (14.2,-2.4) {Descartada};

  \draw[arrow] (conf) -- node[lab,above]{gerar sinal} (sinal);
  \draw[arrow] (sinal) -- node[lab,above]{calcular FFT} (fft);
  \draw[arrow] (fft) -- node[lab,above]{avaliar} (aval);
  \draw[arrow] (aval.east) -- node[lab,above,pos=0.4]{regra de ouro / fora de $R^3$ / primeira leitura} (pers.west);
  \draw[arrow] (aval.south) to[out=-90,in=180] node[lab,left,pos=0.25]{todos dentro de $R^3$} (desc.west);
\end{tikzpicture}

\caption{Máquina de estados do processamento de uma leitura simulada.}
\label{fig:estados}
\caption*{\footnotesize Fonte: elaboração própria, conforme arquitetura do projeto.}
\end{figure}

\section{Requisitos e rastreabilidade}

A especificação formaliza 19 requisitos funcionais (FR-001 a FR-019) e 7 não
funcionais (NFR-001 a NFR-007), detalhados em `SPECIFICATION.md`. A
Tabela~\ref{tab:requisitos} agrupa os requisitos por módulo de implementação.

\begin{table}[htbp]
\centering
\caption{Agrupamento dos requisitos por módulo implementado.}
\label{tab:requisitos}
\begin{tabular}{@{}p{4.4cm}p{8.6cm}@{}}
\toprule
\textbf{Módulo} & \textbf{Requisitos associados} \\
\midrule
`signal_generator` & FR-001, FR-002, FR-003 \\
`fft_processor` & FR-004, FR-005, FR-006, FR-019, NFR-001 \\
`discard_engine` & FR-007, FR-008, FR-009, FR-010, FR-011 \\
`persistence` & FR-011, FR-012, FR-013, NFR-003, NFR-004 \\
`api` & FR-016, FR-017, FR-018, NFR-002 \\
`frontend` & FR-014, FR-015, FR-016, FR-017, FR-018, FR-019 \\
\bottomrule
\end{tabular}
\caption*{\footnotesize Fonte: SPECIFICATION.md do projeto.}
\end{table}

\section{Geração de sinais sintéticos}

O módulo `signal_generator` implementa a síntese no domínio do tempo. Para cada
defeito do catálogo, um conjunto de \emph{componentes espectrais} define a
ordem (múltiplo de \(f_r\)), o peso relativo e, quando aplicável, a geração de
bandas laterais em \(\pm 1X\). O sinal é composto pela soma de senoides
\begin{equation}
\label{eq:sintese}
  s(t) = \sum_i A_i \sin(2\pi f_i t + \varphi_i) + n(t),
\end{equation}
em que \(f_i = ordem_i \cdot f_r\), \(A_i = peso_i \cdot severidade\) e
\(n(t)\) é ruído branco gaussiano com desvio padrão igual ao parâmetro de ruído
de fundo. A Tabela~\ref{tab:perfis} resume os perfis implementados
(Tabela~\ref{tab:catalogo} contém o catálogo completo).

\begin{table}[htbp]
\centering
\caption{Perfis espectrais implementados (ordens e pesos relativos) para defeitos representativos.}
\label{tab:perfis}
\begin{tabular}{@{}p{4.2cm}p{8.8cm}@{}}
\toprule
\textbf{Tipo} & \textbf{Componentes: ordem (peso)} \\
\midrule
sem\_defeito & 1X (0,20); 2X (0,03); 3X (0,015) \\
desbalanceamento & 1X (1,00); 2X (0,07); 3X (0,04) \\
desalinhamento\_angular & 1X (0,50); 2X (0,80); 3X (0,25); 1,5X (0,15); 2,5X (0,12); 4X (0,12) \\
desalinhamento\_paralelo & 1X (0,40); 2X (1,00); 4X (0,40); 6X (0,18); 8X (0,08) \\
rocamento & 1X (1,00); 2X (0,50); 3X (0,30); 4X (0,20); 5X (0,12); + sub-harmônicos se severo \\
mancal\_frouxo & 1X--10X com pesos 0,70 a 0,08; + fracionários 0,5X/1,5X/2,5X \\
acoplamento\_defeituoso & 1X (0,55); 2X (0,75); 3X (0,30); 4X (0,15) \\
rolamento\_bpfi & BPFI (1,00); 2$\times$BPFI (0,40); 3$\times$BPFI (0,20) $\pm$ sidebands 1X \\
\bottomrule
\end{tabular}
\caption*{\footnotesize Fonte: `DEFECT_PROFILES` e `_componentes_do_defeito` em `backend/app/domain/signal_generator.py`.}
\end{table}

Duas decisões de engenharia são relevantes. A primeira é a \emph{derivação da
frequência a partir da rotação}: todas as frequências são recalculadas como
\(f = ordem \cdot RPM/60\), nunca armazenadas como valores absolutos fixos, o
que mantém a coerência cinemática em qualquer regime. A segunda é o
\emph{determinismo da fase por ponto de medição}: a fase de cada componente é
gerada a partir de uma semente derivada do identificador do ponto (referência
de keyphasor estável), o que torna as leituras sucessivas de um mesmo ponto
comparáveis em fase --- condição necessária para a verificação \(R^3\) ---
enquanto o ruído permanece estocástico. O trecho a seguir ilustra o núcleo da
síntese:

\begin{lstlisting}[language=Python, caption={Núcleo da síntese de sinal em signal_generator.py.}, label={lst:gen}]
for componente in _componentes_do_defeito(tipo_defeito, severidade, rng_fase):
    freq = componente.ordem * freq_rotacao_hz      # ordem x (RPM/60)
    amplitude = componente.peso * severidade
    fase = rng_fase.uniform(0, 2*pi)
    sinal += amplitude * np.sin(2*pi*freq*t + fase)
    if componente.sidebands_1x:
        for f_lat in (freq - fr, freq + fr):       # bandas laterais em +-1X
            if f_lat > 0:
                sinal += 0.2*amplitude * np.sin(2*pi*f_lat*t + fase_lat)
sinal += rng_ruido.normal(0, ruido_efetivo, n)
\end{lstlisting}

Para os defeitos com razão subsíncrona dinâmica (\emph{oil whirl} e \emph{whirl
por atrito}), a razão é amostrada em tempo de geração dentro de uma faixa
especificada (\(0{,}39\)--\(0{,}48\) e \(0{,}30\)--\(0{,}70\), respectivamente),
reproduzindo a variabilidade desses fenômenos. Para defeitos não lineares
severo-dependentes (roçamento com severidade acima de \(0{,}6\) e mancal frouxo),
sub-harmônicos e fracionários são adicionados, e o piso de ruído é elevado
(roçamento e folga truncam a forma de onda), conforme a literatura
\cite{randall2011,scheffer2004}.

\section{Processamento espectral (FFT)}

O módulo `fft_processor` calcula a FFT real (via `numpy.fft.rfft`), extrai os
picos e as métricas. A estimativa de \(F_{max}\) de cada defeito é feita pela
maior ordem do perfil (com piso de ordem 10) multiplicada por uma margem de
bandas laterais de \(2{,}0\):
\begin{equation}
\label{eq:fmax}
  F_{max} = \max\left(ordem_{maxima}, 10\right)\times 2{,}0 \times f_r .
\end{equation}
A validação de Nyquist rejeita configurações com \(f_s \le 2F_{max}\) antes de
qualquer processamento (NFR-001). A detecção de picos considera máximos locais
da magnitude acima de um \emph{limiar} absoluto, definido como o limiar relativo
configurável multiplicado pela maior amplitude do espectro:
\begin{equation}
\label{eq:limiar}
  \lambda = \lambda_{rel} \cdot \max_k |X[k]|,
\end{equation}
com os picos ordenados por amplitude e limitados a um máximo configurável. Para
cada pico, registram-se frequência, amplitude (escala de pico, \(2|X|/N\)) e
fase (em graus), compondo os vetores \(R^3\). As métricas RMS total, RMS de
picos, RMS de ruído e valor DC são calculadas conforme a
Seção~\ref{sec:r3} do Capítulo~\ref{cap:fundamentacao}, e o sinal de tempo é
decimado para a visualização.

\section{Motor de descarte}

O módulo `discard_engine` implementa a regra de ouro e o Paralelepípedo de
Descarte (Equações~\eqref{eq:regra-ouro} e~\eqref{eq:descarte} do
Capítulo~\ref{cap:fundamentacao}). Sua interface recebe a leitura atual, a
última leitura efetivamente persistida do ponto (ou `None` na primeira) e as
tolerâncias configuráveis (desvio de rotação, \(\delta_f\), \(\delta_A\) e
\(\delta_\phi\)). O resultado é uma decisão enumerada:

\begin{lstlisting}[language=Python, caption={Decisões possíveis do motor de descarte.}, label={lst:decisoes}]
PERSISTIR_PRIMEIRA_LEITURA   # nao ha referencia (FR-011)
PERSISTIR_REGRA_DE_OURO      # |dRPM| > desvio e ha picos (FR-007)
PERSISTIR_FORA_TOLERANCIA    # algum pico fora do R^3 (FR-010)
DESCARTAR                    # todos os picos dentro do R^3 (FR-009)
\end{lstlisting}

A comparação de cada pico atual é feita contra o pico de referência mais
próximo em frequência; se não houver pico de referência ou se qualquer desvio
exceder a tolerância, a leitura é persistida. Como o motor é uma função pura
(entrada \(\rightarrow\) saída), todas as combinações são testáveis
unitariamente sem banco de dados, e a concorrência sobre o ponteiro de última
leitura é tratada na camada de persistência, dentro da mesma transação que grava
a nova leitura.

\section{Persistência e modelo de dados}

A camada `persistence` utiliza SQLAlchemy 2.x assíncrono e segue a hierarquia
Planta~$\rightarrow$~Área~$\rightarrow$~Máquina~$\rightarrow$~Ponto. O schema
lógico (Tabela~\ref{tab:schema}) materializa a arquitetura não normalizada
adotada (NFR-003): os picos \(R^3\) de cada leitura são armazenados como um
documento JSON em coluna `JSONB`, priorizando a velocidade de ingestão em
relação à conveniência de consulta, conforme as diretrizes de origem.

\begin{table}[htbp]
\centering
\caption{Visão lógica do modelo de dados persistido.}
\label{tab:schema}
\begin{tabular}{@{}p{5.2cm}p{7.8cm}@{}}
\toprule
\textbf{Relação} & \textbf{Conteúdo principal} \\
\midrule
plantas / areas / maquinas & hierarquia com chaves estrangeiras \\
pontos & nó de medição; `ultima\_leitura\_persistida\_id` (ponteiro crítico) \\
leituras\_persistidas & picos\_r3 (JSONB), rms\_total/ruido/picos, valor\_dc, rotacao, níveis, timestamp \\
leituras\_trash & leitura descartada + `motivo\_descarte` \\
snapshots\_defeito & par sensor + tipo de defeito + vínculo com a leitura \\
\bottomrule
\end{tabular}
\caption*{\footnotesize Fonte: `.specs/codebase/ARCHITECTURE.md` e migração Alembic do projeto.}
\end{table}

A primeira migração Alembic cria as oito relações da
Tabela~\ref{tab:schema}. O estado compartilhado crítico é o ponteiro
`pontos.ultima\_leitura\_persistida\_id`: a decisão de descarte lê esse ponteiro
e grava a nova leitura na \emph{mesma transação}, com lock de linha, evitando
corridas entre simulações concorrentes do mesmo ponto. Leituras aprovadas são
escritas em `leituras_persistidas` e atualizam o ponteiro; leituras descartadas
vão para `leituras_trash`, sem alterar a referência.

\section{API REST}

A API é exposta pelo FastAPI com validação de schema Pydantic v2. Os endpoints
implementados são resumidos na Tabela~\ref{tab:endpoints}; os contratos
detalhados (requisição e resposta) constam do Apêndice~\ref{ap:contratos}.

\begin{table}[htbp]
\centering
\caption{Endpoints REST implementados.}
\label{tab:endpoints}
\begin{tabular}{@{}p{4.6cm}p{4.4cm}p{4cm}@{}}
\toprule
\textbf{Endpoint} & \textbf{Descrição} & \textbf{Requisitos} \\
\midrule
POST /simulacoes & gera sinal, FFT e decide descarte & FR-001 a FR-019, NFR-001 \\
POST /snapshots & registra par sensor+defeito & FR-016 \\
GET /health & verificação de saúde do serviço & --- \\
\bottomrule
\end{tabular}
\caption*{\footnotesize Fonte: `backend/app/api/` do projeto.}
\end{table}

A resposta de `POST /simulacoes` contém `sinal_tempo` (decimado), `picos_r3`,
`rms_total`, `rms_ruido`, `rms_picos`, `valor_dc`, `rotacao`, `limiar_picos` e
`limiar_amplitude`, a decisão (`descartada` e `motivo_descarte`),
`tempo_processamento_ms` e `taxa_descarte_acumulada`. Erros são mapeados como
`422` (entrada inválida ou violação de Nyquist), `404` (ponto inexistente) e
`500` (falha interna de persistência, sem corromper o ponteiro de referência).

\section{Front-end}

O front-end é uma SPA em React com gráficos Recharts, organizada em
componentes: `PainelSimulacao` (formulário), `PainelSinal` (sinal e RMS),
`GraficoFFT` (barras de espectro), `IndicadoresSimulacao` (tempo e taxa de
descarte), `ChecklistValidacao`, `SnapshotDefeito` e o componente `Help`
(tooltips de ajuda). O formulário reúne os parâmetros de máquina e de
aquisição, incluindo um controle de limiar relativo de picos; a validação no
cliente espelha o critério de Nyquist e bloqueia a submissão de parâmetros
inválidos (FR-018, NFR-001).

O gráfico de FFT exibe cada barra com a frequência em hertz e a ordem
normalizada \(N\) (frequência dividida pela rotação em hertz), uma linha
horizontal com o limiar absoluto e, abaixo do gráfico, o valor da rotação em
hertz (FR-015). O registro de \emph{snapshot} associa o sensor ao tipo de
defeito simulado e persiste via `POST /snapshots` (FR-016). O client HTTP
tipado trata erros de rede e erros de API de forma explícita, exibindo
mensagens compreensíveis ao usuário.

\section{Infraestrutura e execução}

O ambiente de execução local é definido por Docker Compose com três serviços:
`db` (PostgreSQL 16 com volume persistente), `backend` (imagem Python que executa
`alembic upgrade head` e inicia o Uvicorn na porta 8000) e `frontend` (build de
produção servido por nginx na porta 80, com proxy de `/api/` para o back-end).
Os scripts `start.sh` e `stop.sh` automatizam, respectivamente, o setup (checagem
do Docker, cópia do `.env`, build de imagem se ausente, subida com espera de
health e semeadura de um ponto de demonstração) e a derrubada do ambiente
(inclusive com remoção opcional de volumes).

\section{Decisões de engenharia (ADRs)}

As principais decisões de arquitetura, registradas no projeto como ADRs, foram:

\begin{itemize}
  \item \textbf{JSONB para picos \(R^3\)} (NFR-003): banco relacional com colunas
  JSON para os vetores de picos, priorizando ingestão; consultas por pico
  individual exigem operadores JSON do PostgreSQL (reavaliar com indexação se
  necessário).
  \item \textbf{Processamento síncrono no MVP}: chamada request/response direta
  via FastAPI; sem fila de mensagens. Reavaliar se o volume concorrente (NFR-005)
  crescer.
  \item \textbf{Sem autenticação nesta fase} (NFR-006): uso local/interno, CORS
  restrito ao front de desenvolvimento; a validação Pydantic é a única barreira
  no MVP.
  \item \textbf{Determinismo de fase por ponto} e \textbf{limiar relativo
  configurável}: decisões de domínio que habilitam, respectivamente, a
  comparabilidade de fase no \(R^3\) e o descarte de ruído de fundo (FR-019).
\end{itemize}

\section{Observabilidade}

Os logs de simulação são estruturados por requisição (parâmetros, decisão de
descarte e tempo de processamento). A resposta de cada simulação inclui o tempo
de processamento (FR-017) e a taxa de descarte acumulada (NFR-004), calculada
pela camada de persistência como a proporção de leituras descartadas em relação
às avaliadas por ponto.
